\documentclass[12pt, a4paper]{report}

% ── Packages ──────────────────────────────────────────────────────────────────
% NOTE: Compile with XeLaTeX (required for Unicode Arabic text).
%   Overleaf: Menu → Compiler → XeLaTeX
%   Local:    xelatex Rafeeq_Speech_Recognition_Documentation.tex
\usepackage[a4paper, margin=2.5cm]{geometry}
\usepackage{fontspec}
  \setmainfont{Latin Modern Roman}
\usepackage{polyglossia}
  \setmainlanguage{english}
  \setotherlanguage{arabic}
  \newfontfamily\arabicfont[Script=Arabic, Scale=1.1]{Amiri}
% \textarabic{} is provided automatically by polyglossia (text<lang>{} pattern)
\usepackage{microtype}
\usepackage{graphicx}
\usepackage[table]{xcolor}
\usepackage{booktabs}
\usepackage{tabularx}
\usepackage{longtable}
\usepackage{array}
\usepackage{multirow}
\usepackage{listings}
\usepackage{fancyhdr}
\usepackage{titlesec}
\usepackage{titling}
\usepackage{hyperref}
\usepackage{enumitem}
\usepackage{parskip}
\usepackage{setspace}
\usepackage{tocloft}
\usepackage{mdframed}
\usepackage{amsmath}
\usepackage{caption}

% ── Colours ───────────────────────────────────────────────────────────────────
\definecolor{navyblue}{HTML}{1F3864}
\definecolor{steelblue}{HTML}{2E75B6}
\definecolor{lightblue}{HTML}{EBF3FB}
\definecolor{codegray}{HTML}{F2F2F2}
\definecolor{darkgray}{HTML}{404040}
\definecolor{tableheader}{HTML}{1F3864}
\definecolor{rowalt}{HTML}{EBF3FB}

% ── Hyperref setup ────────────────────────────────────────────────────────────
\hypersetup{
    colorlinks=true,
    linkcolor=steelblue,
    urlcolor=steelblue,
    citecolor=steelblue,
    pdftitle={Rafeeq Speech Recognition Technical Documentation},
    pdfauthor={Rafeeq Team — Zewail City},
}

% ── Section title formatting ──────────────────────────────────────────────────
\titleformat{\chapter}[block]
  {\normalfont\Large\bfseries\color{navyblue}}
  {\thechapter.}{1em}{}
  [\vspace{2pt}\color{steelblue}\titlerule]

\titleformat{\section}
  {\normalfont\large\bfseries\color{navyblue}}
  {\thesection}{1em}{}

\titleformat{\subsection}
  {\normalfont\normalsize\bfseries\color{steelblue}}
  {\thesubsection}{1em}{}

\titleformat{\subsubsection}
  {\normalfont\normalsize\bfseries\color{darkgray}}
  {\thesubsubsection}{1em}{}

\titlespacing*{\chapter}{0pt}{10pt}{20pt}
\titlespacing*{\section}{0pt}{16pt}{8pt}
\titlespacing*{\subsection}{0pt}{12pt}{6pt}

% ── Line spacing ──────────────────────────────────────────────────────────────
\setstretch{1.15}

% ── Code listing style ────────────────────────────────────────────────────────
\lstdefinestyle{rafeeqcode}{
    backgroundcolor=\color{codegray},
    basicstyle=\ttfamily\small,
    breakatwhitespace=false,
    breaklines=true,
    captionpos=b,
    frame=single,
    framesep=4pt,
    rulecolor=\color{steelblue},
    xleftmargin=12pt,
    xrightmargin=4pt,
    aboveskip=8pt,
    belowskip=8pt,
    showstringspaces=false,
    tabsize=2,
    keywordstyle=\bfseries\color{navyblue},
    commentstyle=\itshape\color{darkgray},
    stringstyle=\color{steelblue},
}
\lstset{style=rafeeqcode}

% ── Table helpers ─────────────────────────────────────────────────────────────
\newcommand{\thead}[1]{%
  \cellcolor{tableheader}\textcolor{white}{\textbf{#1}}}
\newcommand{\trow}[1]{\rowcolor{rowalt}#1}
\renewcommand{\arraystretch}{1.35}

% ── Page header / footer ──────────────────────────────────────────────────────
\pagestyle{fancy}
\fancyhf{}
\fancyhead[L]{\small\color{darkgray}\textit{Rafeeq Smart Wheelchair}}
\fancyhead[R]{\small\color{darkgray}\textit{Speech Recognition — Technical Documentation}}
\fancyfoot[C]{\small\thepage}
\renewcommand{\headrulewidth}{0.4pt}
\renewcommand{\headrule}{\hbox to\headwidth{\color{steelblue}\leaders\hrule height\headrulewidth\hfill}}

% ── TOC formatting ────────────────────────────────────────────────────────────
\renewcommand{\cftchapfont}{\bfseries\color{navyblue}}
\renewcommand{\cftsecfont}{\color{darkgray}}
\setlength{\cftbeforechapskip}{6pt}

% ─────────────────────────────────────────────────────────────────────────────
\begin{document}

% ── COVER PAGE ────────────────────────────────────────────────────────────────
\begin{titlepage}
  \centering
  \vspace*{2cm}

  {\fontsize{28}{34}\selectfont\bfseries\color{navyblue}
    Rafeeq Smart Wheelchair\par}

  \vspace{0.4cm}
  {\color{steelblue}\rule{0.8\textwidth}{2pt}\par}
  \vspace{0.4cm}

  {\Large\bfseries\color{darkgray}
    Speech Recognition Subsystem\par
    Technical Documentation\par}

  \vspace{0.6cm}
  {\normalsize\itshape\color{darkgray}
    Offline Arabic Voice Control: Design, Evaluation, and Architecture\par}

  \vspace{2.5cm}

  \begin{tabular}{rl}
    \textbf{Project}      & Rafeeq (SmartWheel) \\[4pt]
    \textbf{Institution}  & Zewail City of Science and Technology \\[4pt]
    \textbf{Date}         & April 24, 2026 \\[4pt]
    \textbf{Version}      & v1.0 \\
  \end{tabular}

  \vfill
  {\color{steelblue}\rule{0.8\textwidth}{1pt}\par}
  \vspace{0.3cm}
  {\small\color{darkgray}
    Submitted to: Dr.\ Moustafa Elshafei \& Dr.\ Ahmed Abd-Rabou\par}
\end{titlepage}

% ── TABLE OF CONTENTS ─────────────────────────────────────────────────────────
\tableofcontents
\newpage

% ═════════════════════════════════════════════════════════════════════════════
\chapter{Overview}
% ═════════════════════════════════════════════════════════════════════════════

\section{Purpose of This Document}

This document provides a comprehensive technical reference for the speech recognition
subsystem of the Rafeeq smart wheelchair. It covers the three architectures evaluated
during development, the quantitative comparison that led to the final design decision,
and a detailed breakdown of the selected architecture --- the local TFLite KWS model ---
including its signal processing pipeline, neural network design, inference engine, and
integration with the wheelchair control system.

\section{System Context}

The speech recognition subsystem converts the wheelchair user's spoken Egyptian Arabic
commands into one of 11 discrete motion and navigation intents. It must operate:

\begin{itemize}[noitemsep]
  \item Completely offline (no internet connection at any time)
  \item On a Raspberry Pi~5 with shared CPU resources (navigation, health monitoring,
        and motor control also running concurrently)
  \item Within safety-critical latency constraints ($<$100\,ms for motion commands)
  \item With high reliability for short, single-word Egyptian dialect utterances
\end{itemize}

The 11 supported command intents are listed in Table~\ref{tab:intents}.

\begin{table}[h!]
\centering
\caption{Supported wheelchair command intents}
\label{tab:intents}
\begin{tabular}{clc}
\toprule
\thead{Index} & \thead{Intent Label} & \thead{Action} \\
\midrule
\trow{1}  & \texttt{MOVE\_FORWARD}    & Move wheelchair forward \\
2         & \texttt{MOVE\_BACKWARD}   & Move wheelchair backward \\
\trow{3}  & \texttt{TURN\_LEFT}       & Turn left \\
4         & \texttt{TURN\_RIGHT}      & Turn right \\
\trow{5}  & \texttt{STOP}             & Emergency stop \\
6         & \texttt{GO\_TO\_KITCHEN}  & Navigate to kitchen \\
\trow{7}  & \texttt{GO\_TO\_BATHROOM} & Navigate to bathroom \\
8         & \texttt{GO\_TO\_BEDROOM}  & Navigate to bedroom \\
\trow{9}  & \texttt{GO\_TO\_LIVING\_ROOM} & Navigate to living room \\
10        & \texttt{SPEED\_UP}        & Increase speed \\
\trow{11} & \texttt{SLOW\_DOWN}       & Decrease speed \\
\bottomrule
\end{tabular}
\end{table}

% ═════════════════════════════════════════════════════════════════════════════
\chapter{Three Architectures Evaluated}
% ═════════════════════════════════════════════════════════════════════════════

Three fundamentally different approaches to offline Arabic speech recognition were
built, deployed, and benchmarked on the target hardware (Raspberry Pi~5,
ARM Cortex-A76, Ubuntu 24.04 LTS).

% ─────────────────────────────────────────────────────────────────────────────
\section{Approach A --- Local TFLite KWS Model (MFCC + CNN)}
% ─────────────────────────────────────────────────────────────────────────────

A custom-trained Keyword Spotting (KWS) model based on Mel-Frequency Cepstral
Coefficients (MFCC) feature extraction followed by a Convolutional Neural Network (CNN).
The model is trained exclusively on the 11 wheelchair commands in Egyptian Arabic,
exported as a TFLite binary, and runs in the \texttt{ai\_edge\_litert} inference engine.

\textbf{Processing pipeline:}
\begin{enumerate}[noitemsep]
  \item Raw 16-bit PCM audio is captured in a 1.5-second sliding window at 16\,kHz
  \item Audio is converted to a \texttt{float32} array and normalised to $[-1, 1]$
  \item \texttt{librosa} extracts 13 MFCC coefficients across 47 time frames
        $\rightarrow$ tensor shape \texttt{[1, 13, 47, 1]}
  \item CNN model runs inference on the MFCC tensor
  \item Softmax output over 11 classes $\rightarrow$ argmax gives the predicted command
  \item A confidence threshold ($\geq 80\%$) must be met for the command to be accepted
\end{enumerate}

\begin{mdframed}[backgroundcolor=lightblue, linecolor=steelblue, linewidth=1pt]
\textbf{Key metrics:} \quad
Model size: \textbf{$\sim$100\,KB} \quad|\quad
Latency: \textbf{$<$10\,ms} \quad|\quad
CPU: \textbf{$<$2\%}
\end{mdframed}

% ─────────────────────────────────────────────────────────────────────────────
\section{Approach B --- Vosk (Kaldi TDNN, Offline ASR)}
% ─────────────────────────────────────────────────────────────────────────────

Vosk is an offline speech recognition library based on Kaldi's Time-Delay Neural
Network (TDNN) acoustic model. It performs full speech-to-text transcription; a keyword
matcher then maps the transcript to a command intent.

\textbf{Processing pipeline:}
\begin{enumerate}[noitemsep]
  \item Continuous audio stream fed to the Vosk recogniser
  \item Kaldi TDNN decodes audio into a text transcript (e.g.\ \textarabic{اطلع قدام})
  \item Keyword matching maps the transcript to one of 11 intents
\end{enumerate}

\begin{mdframed}[backgroundcolor=lightblue, linecolor=steelblue, linewidth=1pt]
\textbf{Key metrics:} \quad
Model size: \textbf{$\sim$300\,MB} \quad|\quad
Latency: \textbf{$\sim$120\,ms} \quad|\quad
CPU: \textbf{35--45\%} \quad|\quad
Accuracy: \textbf{94\%}
\end{mdframed}

% ─────────────────────────────────────────────────────────────────────────────
\section{Approach C --- Hybrid (Whisper.cpp + BoW Intent Classifier)}
% ─────────────────────────────────────────────────────────────────────────────

A 3-layer hybrid pipeline combining a TinyML wake-word detector, Whisper.cpp (C++ port
of OpenAI Whisper) for Arabic speech-to-text, and a Bag-of-Words (BoW) + Fully
Connected TFLite classifier for intent mapping.

\textbf{Layer breakdown:}
\begin{itemize}[noitemsep]
  \item \textbf{Layer 1} --- Wake-word detector (TFLite CNN): always-on, detects
        \textit{``rafeeq''} to activate the pipeline
  \item \textbf{Layer 2} --- Whisper.cpp (\texttt{ggml-base.bin}): transcribes the
        3-second command window to raw Arabic text
  \item \textbf{Layer 3} --- BoW + FC TFLite classifier: maps Whisper transcript to
        one of 11 intents
\end{itemize}

\begin{mdframed}[backgroundcolor=lightblue, linecolor=steelblue, linewidth=1pt]
\textbf{Key metrics:} \quad
Total size: \textbf{$\sim$210\,MB} \quad|\quad
Latency: \textbf{2--3\,s} \quad|\quad
CPU: \textbf{$\sim$80\%} \quad|\quad
Accuracy: \textbf{Below target}
\end{mdframed}

% ═════════════════════════════════════════════════════════════════════════════
\chapter{Comparative Analysis and Decision}
% ═════════════════════════════════════════════════════════════════════════════

\section{Benchmarking Table}

\begin{table}[h!]
\centering
\caption{Head-to-head comparison of all three architectures}
\label{tab:comparison}
\begin{tabularx}{\textwidth}{>{\bfseries}lXXX}
\toprule
\thead{Criterion} & \thead{Local KWS} & \thead{Vosk (Kaldi)} & \thead{Hybrid (Whisper)} \\
\midrule
\trow{Model Size}        & $\sim$100 KB   & $\sim$300 MB   & $\sim$210 MB \\
Latency                  & $<$10 ms       & $\sim$120 ms   & 2,000--3,000 ms \\
\trow{CPU Usage}         & $<$2\%         & 35--45\%       & $\sim$80\% \\
Offline Operation        & Yes            & Yes            & Yes \\
\trow{Short-cmd Accuracy}& High (direct)  & 94\%           & Below target \\
Egyptian Dialect         & Custom-trained & Limited        & Good (general) \\
\trow{Customisation}     & Full control   & Low            & Medium \\
RAM Requirement          & $<$5 MB        & $\sim$300 MB   & $\sim$250 MB \\
\trow{Safety Suitability}& \textbf{Excellent} & Acceptable & \textcolor{red}{\textbf{Unacceptable}} \\
\bottomrule
\end{tabularx}
\end{table}

\section{Why the Local KWS Model Was Selected}

Three factors drove the final decision:

\subsection{Factor 1 --- Latency (Safety-Critical)}

The wheelchair is a safety-critical system. A \texttt{STOP} command issued in an
emergency must be acted upon in milliseconds, not seconds. The Hybrid Whisper approach
requires 2--3 seconds --- an unacceptable delay. The local KWS model responds in under
10\,ms, which is:

\begin{center}
\textbf{200$\times$ faster than Vosk \quad|\quad 300$\times$ faster than Hybrid Whisper}
\end{center}

\subsection{Factor 2 --- Resource Efficiency}

The Raspberry Pi~5 runs four concurrent subsystems: navigation (SLAM + Nav2), health
monitoring, motor control (serial UART), and speech recognition. Vosk's 35--45\% CPU
usage and 300\,MB RAM footprint directly compete with the navigation stack. The local
KWS model consumes less than 2\% CPU and only 100\,KB of storage ---
\textit{effectively invisible to the rest of the system.}

\subsection{Factor 3 --- Customisation for Egyptian Arabic}

Both Vosk and Whisper are general-purpose ASR systems not optimised for short,
single-word Egyptian colloquial commands. The local KWS model is trained
\textbf{exclusively} on the exact 11 commands the wheelchair needs, in the exact
dialect the patient uses. This domain-specific training gives higher reliability than
any general-purpose ASR system.

% ═════════════════════════════════════════════════════════════════════════════
\chapter{Local KWS Model: Detailed Architecture}
% ═════════════════════════════════════════════════════════════════════════════

\section{Full Pipeline Overview}

\begin{lstlisting}[caption={Complete voice-to-command pipeline},
                   label={lst:pipeline}]
[Microphone - USB, 16 kHz, 16-bit PCM]
        |
        v
[VAD Gate - RMS amplitude threshold]
  RMS < 50  -->  discard chunk, keep listening
  RMS >= 50 -->  start filling KWS window
        |
        v
[1.5-second Audio Buffer - 24,000 samples @ 16 kHz]
        |
        v
[Audio Preprocessing]
  1. Convert int16 bytes --> float32, normalise to [-1, 1]
  2. librosa.effects.trim(top_db=20) -- strip silence
  3. Pad or truncate to exactly 24,000 samples
        |
        v
[MFCC Feature Extraction - librosa]
  n_mfcc    = 13 coefficients
  n_fft     = 512  (32 ms window)
  hop_length = 512
  Output: (13, 47) --> reshape to [1, 13, 47, 1]
        |
        v
[TFLite CNN Inference - rafeeq_model.tflite (~100 KB)]
  Input:  [1, 13, 47, 1]  float32
  Output: [1, 11]         float32 softmax probabilities
        |
        v
[Confidence Gate - threshold >= 0.80]
  confidence < 0.80  -->  UNKNOWN, return to idle
  confidence >= 0.80 -->  accept command
        |
        v
[Command Dispatch - 11 intents]
  --> Serial UART to Arduino / Nav2 ROS2 goal
\end{lstlisting}

% ─────────────────────────────────────────────────────────────────────────────
\section{Audio Capture and VAD Gate}
% ─────────────────────────────────────────────────────────────────────────────

The \texttt{AudioRecorder} class maintains a persistent PyAudio stream at 16{,}000\,Hz
mono 16-bit PCM, read in chunks of 512 frames ($\approx$32\,ms each).

A Voice Activity Detection (VAD) gate based on RMS amplitude filters out silence before
any DSP or ML work is performed:

\begin{equation}
  \mathrm{RMS} = \sqrt{\frac{1}{N} \sum_{i=1}^{N} x_i^2}
\end{equation}

If $\mathrm{RMS} < \theta_{\mathrm{VAD}}$ (default: $\theta = 50$), the chunk is
discarded and the loop iterates. This keeps CPU usage near zero during silence.

When a volume spike is detected, a full 1.5-second window (47 chunks of 512 frames)
is captured into a buffer before MFCC extraction begins.

% ─────────────────────────────────────────────────────────────────────────────
\section{MFCC Feature Extraction}
% ─────────────────────────────────────────────────────────────────────────────

MFCC (Mel-Frequency Cepstral Coefficients) approximates how the human auditory system
perceives sound by applying a mel-scale filterbank to the power spectrum. Extraction
steps:

\begin{enumerate}[noitemsep]
  \item Trim leading/trailing silence: \texttt{librosa.effects.trim(top\_db=20)}
  \item Pad or truncate audio to exactly 24{,}000 samples (1.5\,s at 16\,kHz)
  \item Compute 13 MFCC coefficients:
        \begin{itemize}[noitemsep]
          \item \texttt{n\_fft = 512} (32\,ms analysis window)
          \item \texttt{hop\_length = $\lfloor 16000 \times 1.5 / 47 \rfloor$}
                $\rightarrow$ exactly 47 time frames
        \end{itemize}
  \item Pad/truncate the time axis to exactly 47 frames
  \item Reshape to \texttt{[1, 13, 47, 1]} --- the model's expected input tensor
\end{enumerate}

The resulting tensor encodes the spectral envelope across 13 perceptual frequency bands
over 47 time steps, covering the full 1.5-second command window.

% ─────────────────────────────────────────────────────────────────────────────
\section{Neural Network Architecture (CNN)}
% ─────────────────────────────────────────────────────────────────────────────

The \texttt{rafeeq\_model.tflite} CNN processes the $[13 \times 47 \times 1]$ input as
a 2-D image where the height axis represents frequency bands and the width axis
represents time.

\begin{lstlisting}[caption={KWS CNN architecture summary}]
Input:  [1, 13, 47, 1]
  --> Conv2D   (32 filters, 3x3, ReLU, padding=same)
  --> MaxPool2D (2x2)
  --> Conv2D   (64 filters, 3x3, ReLU, padding=same)
  --> MaxPool2D (2x2)
  --> Flatten
  --> Dense    (128 units, ReLU)
  --> Dropout  (0.25)
  --> Dense    (11 units, Softmax)
Output: [1, 11]  -- probability over 11 command classes
\end{lstlisting}

The model is exported with dynamic-range quantisation, reducing the binary to
$\sim$100\,KB while preserving inference accuracy.

% ─────────────────────────────────────────────────────────────────────────────
\section{TFLite Inference Engine}
% ─────────────────────────────────────────────────────────────────────────────

The inference engine uses \texttt{ai\_edge\_litert} (Google's official replacement for
the discontinued \texttt{tflite\_runtime}), providing:

\begin{itemize}[noitemsep]
  \item XNNPACK delegate for ARM NEON SIMD optimised operations on the Cortex-A76
  \item Support for \texttt{FULLY\_CONNECTED} op version 12+ (unavailable in old
        \texttt{tflite\_runtime})
  \item Persistent interpreter with pre-allocated tensors (zero allocation overhead
        per inference)
\end{itemize}

\begin{lstlisting}[language=Python, caption={TFLite inference flow}]
interpreter = TFLiteInterpreter(model_path="rafeeq_model.tflite")
interpreter.allocate_tensors()
in_idx  = interpreter.get_input_details()[0]["index"]
out_idx = interpreter.get_output_details()[0]["index"]

interpreter.set_tensor(in_idx, mfcc_tensor)  # [1, 13, 47, 1]
interpreter.invoke()                          # < 10 ms on RPi5
probs      = interpreter.get_tensor(out_idx)[0]  # shape: (11,)
best_idx   = int(np.argmax(probs))
confidence = float(probs[best_idx])
\end{lstlisting}

% ─────────────────────────────────────────────────────────────────────────────
\section{Confidence Threshold and Safety Gate}
% ─────────────────────────────────────────────────────────────────────────────

\begin{equation}
  \text{command} =
  \begin{cases}
    \text{INTENT\_LABELS}[\text{argmax}(\mathbf{p})] & \text{if } \max(\mathbf{p}) \geq 0.80 \\
    \text{UNKNOWN (ignored)}                          & \text{if } \max(\mathbf{p}) < 0.80
  \end{cases}
\end{equation}

The 80\% threshold is a critical safety requirement. If the model is uncertain due to
background noise, an unrecognised word, or ambiguous audio, it \textbf{must not} issue
a movement command. It is safer to ask the user to repeat a command than to move in an
unintended direction.

% ─────────────────────────────────────────────────────────────────────────────
\section{Label Set and Class Mapping}
% ─────────────────────────────────────────────────────────────────────────────

\begin{table}[h!]
\centering
\caption{KWS model output class mapping (loaded from \texttt{labels.txt})}
\label{tab:labels}
\begin{tabular}{clll}
\toprule
\thead{Index} & \thead{Label (labels.txt)} & \thead{Dispatch Key} & \thead{Action} \\
\midrule
\trow{0}  & \texttt{go\_to\_bathroom}   & \texttt{GO\_TO\_BATHROOM}   & Navigate to bathroom \\
1         & \texttt{go\_to\_bedroom}    & \texttt{GO\_TO\_BEDROOM}    & Navigate to bedroom \\
\trow{2}  & \texttt{go\_to\_kitchen}    & \texttt{GO\_TO\_KITCHEN}    & Navigate to kitchen \\
3         & \texttt{go\_to\_livingroom} & \texttt{GO\_TO\_LIVING\_ROOM} & Navigate to living room \\
\trow{4}  & \texttt{move\_backward}     & \texttt{MOVE\_BACKWARD}     & Move backward \\
5         & \texttt{move\_forward}      & \texttt{MOVE\_FORWARD}      & Move forward \\
\trow{6}  & \texttt{rafeeq}             & (wake-word activation)       & Opens command window \\
7         & \texttt{sleep}              & (standby mode)               & System standby \\
\trow{8}  & \texttt{stop}              & \texttt{STOP}               & \textbf{Emergency stop} \\
9         & \texttt{turn\_left}         & \texttt{TURN\_LEFT}         & Turn left \\
\trow{10} & \texttt{turn\_right}        & \texttt{TURN\_RIGHT}        & Turn right \\
\bottomrule
\end{tabular}
\end{table}

\textbf{Note:} The \texttt{stop} class (index 8) bypasses all further processing and
immediately dispatches the STOP command --- the shortest possible path from voice to
motor halt.

% ─────────────────────────────────────────────────────────────────────────────
\section{Command Dispatcher}
% ─────────────────────────────────────────────────────────────────────────────

\begin{lstlisting}[language=Python, caption={Command dispatch routing}]
handlers = {
    "MOVE_FORWARD"     : _cmd_move_forward,      # Serial UART --> Arduino
    "MOVE_BACKWARD"    : _cmd_move_backward,
    "TURN_LEFT"        : _cmd_turn_left,
    "TURN_RIGHT"       : _cmd_turn_right,
    "STOP"             : _cmd_stop,              # Zero-velocity (safety-critical)
    "GO_TO_KITCHEN"    : lambda: _cmd_navigate("kitchen"),    # Nav2 ROS2 goal
    "GO_TO_BATHROOM"   : lambda: _cmd_navigate("bathroom"),
    "GO_TO_BEDROOM"    : lambda: _cmd_navigate("bedroom"),
    "GO_TO_LIVING_ROOM": lambda: _cmd_navigate("living_room"),
    "SPEED_UP"         : _cmd_speed_up,
    "SLOW_DOWN"        : _cmd_slow_down,
}
\end{lstlisting}

% ═════════════════════════════════════════════════════════════════════════════
\chapter{Layer 3: Intent Classifier (Text-Based --- Hybrid Pipeline)}
% ═════════════════════════════════════════════════════════════════════════════

A secondary text-based intent classifier (\texttt{text\_classifier.tflite}) was
developed for the Hybrid pipeline's Layer 3. Though the Hybrid approach was not
selected as the primary architecture, it is documented here for completeness.

\section{Architecture}

\begin{lstlisting}[caption={Layer 3 BoW + FC classifier pipeline}]
Input:  Arabic text string (Whisper transcript)
   |
   v
normalize_arabic(text)         -- 6-step normalisation
   |
   v
Bag-of-Words vectorisation     -- 120-token vocabulary
   |
   v
TFLite FC inference:
  Dense(128, ReLU) --> Dropout(0.3)
  --> Dense(64, ReLU) --> Dropout(0.2)
  --> Dense(11, Softmax)
   |
   v
Output: intent label + confidence score
\end{lstlisting}

\section{Arabic Text Normalisation Pipeline (6 Steps)}

The \texttt{normalize\_arabic()} function ensures formal MSA and Egyptian colloquial
variants of the same word map to the same vocabulary token:

\begin{table}[h!]
\centering
\caption{Arabic normalisation steps}
\label{tab:normalisation}
\begin{tabularx}{\textwidth}{clX}
\toprule
\thead{Step} & \thead{Operation} & \thead{Example} \\
\midrule
\trow{1} & Remove diacritics (harakat) \& tatweel
         & \textarabic{يمينًا} $\rightarrow$ \textarabic{يمينا} \\
2 & Unify Alef variants (\textarabic{أ إ آ ٱ} $\rightarrow$ \textarabic{ا})
  & \textarabic{للأمام} $\rightarrow$ \textarabic{للامام} \\
\trow{3} & Unify Ya tail (\textarabic{ى} $\rightarrow$ \textarabic{ي})
         & \textarabic{ابقى} $\rightarrow$ \textarabic{ابقي} \\
4 & Unify Hamza seats (\textarabic{ؤ}$\rightarrow$\textarabic{و}, \textarabic{ئ}$\rightarrow$\textarabic{ي})
  & --- \\
\trow{5} & Strip punctuation, numbers, symbols & \textarabic{قدام!} $\rightarrow$ \textarabic{قدام} \\
6 & Strip prefixes (\textarabic{بال لل ال في و}) --- longest first, $\geq$2 chars remain
  & \textarabic{الحمام} $\rightarrow$ \textarabic{حمام} \\
\bottomrule
\end{tabularx}
\end{table}

% ═════════════════════════════════════════════════════════════════════════════
\chapter{Deployment: Dependency Resolution on Raspberry Pi~5}
% ═════════════════════════════════════════════════════════════════════════════

Three dependency conflicts were encountered and resolved during deployment on the
Raspberry Pi~5 (Ubuntu 24.04 LTS, Python 3.10). These are documented for
reproducibility.

\begin{table}[h!]
\centering
\caption{Dependency conflicts and resolutions}
\label{tab:deps}
\begin{tabularx}{\textwidth}{lXX}
\toprule
\thead{Issue} & \thead{Root Cause} & \thead{Resolution} \\
\midrule
\trow{NumPy 2.2.6 crash}
  & \texttt{tflite\_runtime} compiled against NumPy 1.x;
    NumPy 2.0 restructured internal C API (\texttt{\_ARRAY\_API not found})
  & Downgrade: \texttt{pip install "numpy>=1.23,<2.0"} \\
\texttt{FULLY\_CONNECTED} op v12 not found
  & Old \texttt{tflite\_runtime} only supports op $\leq$v11;
    TF 2.14+ generates v12
  & Migrate: \texttt{pip uninstall tflite-runtime \&\& pip install ai-edge-litert} \\
\trow{\texttt{coverage.types} AttributeError}
  & \texttt{numba} (used by librosa) imports \texttt{coverage} at load time;
    system coverage 6.2 missing \texttt{.types} (added in v7.0)
  & Shadow system pkg: \texttt{pip install "coverage>=7.0"} \\
\bottomrule
\end{tabularx}
\end{table}

% ═════════════════════════════════════════════════════════════════════════════
\chapter{Performance Summary}
% ═════════════════════════════════════════════════════════════════════════════

\section{Final Performance on Raspberry Pi~5}

\begin{table}[h!]
\centering
\caption{Local KWS model --- measured performance metrics}
\label{tab:perf}
\begin{tabular}{ll}
\toprule
\thead{Metric} & \thead{Value} \\
\midrule
\trow{Model binary size}           & $\sim$100 KB (\texttt{rafeeq\_model.tflite}) \\
MFCC extraction time               & 8--15 ms (librosa, CPU) \\
\trow{TFLite inference time}       & $<$10 ms (XNNPACK delegate, ARM NEON) \\
Total pipeline latency             & $\sim$30--50 ms end-to-end \\
\trow{CPU usage during inference}  & $<$2\% \\
RAM footprint                      & $<$5 MB \\
\trow{Confidence threshold}        & 80\% \\
Supported commands                 & 11 \\
\trow{Offline operation}           & Yes --- zero network dependency \\
Audio input window                 & 1.5 s (post-VAD) \\
\bottomrule
\end{tabular}
\end{table}

\section{Model Size Comparison}

\begin{table}[h!]
\centering
\caption{Full comparison across all three evaluated architectures}
\label{tab:fullcomp}
\begin{tabular}{lrrr}
\toprule
\thead{Metric} & \thead{Local KWS} & \thead{Vosk} & \thead{Hybrid Whisper} \\
\midrule
\trow{Storage footprint} & \textbf{100 KB}  & 300 MB  & 210 MB \\
Latency                  & \textbf{$<$10 ms}& 120 ms  & 2,000--3,000 ms \\
\trow{CPU usage}         & \textbf{$<$2\%}  & 35--45\%& $\sim$80\% \\
RAM usage                & \textbf{$<$5 MB} & 300 MB  & 250 MB \\
\trow{Short-cmd accuracy}& High             & 94\%    & Below target \\
Dialect customisation    & Full             & Low     & Medium \\
\trow{Safety suitability}& \textbf{Excellent}& Acceptable & \textcolor{red}{Unacceptable} \\
\bottomrule
\end{tabular}
\end{table}

% ═════════════════════════════════════════════════════════════════════════════
\chapter{Files and Configuration Reference}
% ═════════════════════════════════════════════════════════════════════════════

\section{Key Files}

\begin{table}[h!]
\centering
\caption{Project file inventory}
\label{tab:files}
\begin{tabularx}{\textwidth}{llX}
\toprule
\thead{File} & \thead{Size} & \thead{Description} \\
\midrule
\trow{\texttt{rafeeq\_model.tflite}}       & $\sim$100 KB & Layer 1 KWS CNN --- wake-word and command model \\
\texttt{labels.txt}                         & $<$1 KB      & Class index $\rightarrow$ label mapping \\
\trow{\texttt{text\_classifier.tflite}}    & $\sim$60 KB  & Layer 3 BoW+FC intent classifier (Hybrid) \\
\texttt{metadata.json}                      & $\sim$5 KB   & Vocabulary list and label order for text classifier \\
\trow{\texttt{main.py}}                    & $\sim$30 KB  & Full pipeline: recorder, KWS, classifier, dispatcher \\
\texttt{requirements.txt}                   & $<$1 KB      & Pinned Pi~5 dependencies \\
\trow{\texttt{text\_model\_creator.ipynb}} & $\sim$20 KB  & Colab notebook: dataset, training, TFLite export \\
\bottomrule
\end{tabularx}
\end{table}

\section{Key Configuration Constants}

\begin{table}[h!]
\centering
\caption{Runtime configuration constants in \texttt{main.py}}
\label{tab:constants}
\begin{tabularx}{\textwidth}{llX}
\toprule
\thead{Constant} & \thead{Default} & \thead{Description} \\
\midrule
\trow{\texttt{SAMPLE\_RATE}}             & 16{,}000 Hz     & Audio sample rate (must match training) \\
\texttt{KWS\_WINDOW\_SECS}               & 1.5 s           & Audio buffer duration per KWS inference \\
\trow{\texttt{KWS\_CONFIDENCE\_THRESH}}  & 0.80            & Minimum confidence for wake-word acceptance \\
\texttt{INTENT\_CONFIDENCE\_THRESH}      & 0.80            & Minimum confidence for intent acceptance \\
\trow{\texttt{VOLUME\_THRESHOLD}}        & 50              & Minimum RMS to pass the VAD gate \\
\texttt{N\_MFCC}                         & 13              & Number of MFCC coefficients \\
\trow{\texttt{KWS\_INPUT\_SHAPE}}        & [1,13,47,1]     & Model input tensor shape \\
\texttt{COMMAND\_WINDOW\_SECS}           & 3 s             & Post-wake-word recording (Hybrid pipeline) \\
\trow{\texttt{WHISPER\_LANGUAGE}}        & \texttt{"ar"}   & Language code for Whisper transcription \\
\texttt{WHISPER\_MODEL\_SIZE}            & \texttt{"base"} & Whisper model size (tiny/base/small) \\
\bottomrule
\end{tabularx}
\end{table}

% ═════════════════════════════════════════════════════════════════════════════
\chapter{Next Steps}
% ═════════════════════════════════════════════════════════════════════════════

\begin{enumerate}
  \item \textbf{Real-user data collection} --- Record 50+ utterances per command from
        actual Egyptian Arabic speakers to retrain the KWS model on real-world voice patterns.

  \item \textbf{Motor control integration} --- Replace the stub \texttt{dispatch\_command()}
        handlers with real serial UART commands to the Arduino Mega motor controller,
        completing the voice-to-movement chain.

  \item \textbf{Wake-word gate tuning} --- Calibrate \texttt{VOLUME\_THRESHOLD} and
        \texttt{KWS\_CONFIDENCE\_THRESH} for the acoustic environment of the
        June~1 demonstration venue.

  \item \textbf{STOP command latency validation} --- Measure end-to-end latency from
        start of \textarabic{وقف} utterance to motor zero-velocity command,
        targeting $<$150\,ms total.

  \item \textbf{Multi-user robustness testing} --- Validate accuracy across users with
        different voice characteristics (male, female, elderly) and acoustic environments
        (quiet room vs.\ background noise).
\end{enumerate}

% ═════════════════════════════════════════════════════════════════════════════
\chapter{Glossary}
% ═════════════════════════════════════════════════════════════════════════════

\begin{table}[h!]
\centering
\caption{Technical terms and definitions}
\label{tab:glossary}
\begin{tabularx}{\textwidth}{lX}
\toprule
\thead{Term} & \thead{Definition} \\
\midrule
\trow{KWS}     & Keyword Spotting --- detecting a specific word/phrase in a continuous audio stream \\
MFCC           & Mel-Frequency Cepstral Coefficients --- spectral audio features approximating human auditory perception \\
\trow{TFLite}  & TensorFlow Lite --- lightweight ML inference framework for edge devices \\
CNN            & Convolutional Neural Network --- deep learning architecture applying learnable filters to input tensors \\
\trow{BoW}     & Bag-of-Words --- text encoded as a vector of word occurrence counts \\
ASR            & Automatic Speech Recognition --- converting spoken audio to text \\
\trow{VAD}     & Voice Activity Detection --- detecting speech presence before running the full model \\
XNNPACK        & High-performance ARM neural network inference library used by TFLite as a delegate \\
\trow{UART}    & Universal Asynchronous Receiver-Transmitter --- serial protocol between RPi5 and Arduino \\
SLAM           & Simultaneous Localisation and Mapping --- building a map while tracking position \\
\trow{Nav2}    & Navigation 2 --- the ROS2 navigation stack for autonomous path planning \\
EKF            & Extended Kalman Filter --- sensor fusion algorithm combining odometry and IMU data \\
\trow{FOC}     & Field-Oriented Control --- smooth, efficient BLDC motor control technique \\
\bottomrule
\end{tabularx}
\end{table}

% ─────────────────────────────────────────────────────────────────────────────
\vfill
\begin{center}
  {\color{steelblue}\rule{0.6\textwidth}{1pt}}\\[6pt]
  {\small\color{darkgray}
    Rafeeq (SmartWheel) Project \quad|\quad
    Zewail City of Science and Technology \quad|\quad
    April 2026}
\end{center}

\end{document}
